% AUTO-GENERATED by scripts/cpii_simulator/build_appendix.py; do not edit.
% Deterministic source provenance (SHA256):
% source-sha256: data/processed/cpii_simulator/examples/fictional_companies.csv 36c851ef130ca99f7e985116b01b1f50bffa3603801dabcddfe9bae60199a424
% source-sha256: outputs/tables/cpii_simulator/examples/company_ranking.csv fcce7b7945d1d710429a0b7995f6d479f1f14ae7261b9b28e4af0dbf46d9f309
% source-sha256: outputs/tables/cpii_simulator/examples/pair_gap_matrix.csv 7b632e471aa0d90f4587ae4e3ec946767a6b013fc87313bc4d0adeb3993addd4
% source-sha256: outputs/tables/cpii_simulator/examples/pair_gaps.csv 264e5d5feb75e4933460827f33706e46a92e40ae605ff65ced8558f36873126c
% source-sha256: outputs/tables/cpii_simulator/examples/pair_contributions.csv 755afa4162a3b0240adddf973c7136b320d38a2a32d97ad573d1f0f3d23f6e42
% source-sha256: outputs/tables/cpii_simulator/examples/economic_interpretations.csv af728bcf517b486544e41c63ade504407f94099fe739de3f3603b76461b1116c
% source-sha256: outputs/tables/cpii_simulator/examples/golden_equal_gap_summary.csv c989ac21688473bf9184d2d78d8d149b553f96d0c604d7dcf6c9cb216b00eefd
% source-sha256: outputs/tables/cpii_simulator/examples/sensitivity_summary.csv 49091be5635fd9c12acb029464a2b8da882ea62fcea4f56b740907d6c993f819
% source-sha256: data/processed/cpii_simulator/anchor/POOL_SCQ_P01_P99_EMP_RANK_V1.json 274a907a5b383a4c00a57927ccb46c1488e8bfdee9990a1ed87741ca4a769e6a

\chapter{회사 입력형 CPII 비교와 경제적 차이 해석}
\label{app:cpii_company_comparator}

\noindent\fbox{\parbox{0.96\textwidth}{\textbf{Research Beta 경고.} 방법론 설명용 합성자료이며 실제 기업에 대한 추정이 아님. 현재 고정 앵커의 결합 기준표본과 일부 feature의 지원도가 매우 낮으므로, 결과는 연구·전략분석 보조지표로만 사용한다. 차별·역차별·위법성·손해 또는 인과효과를 판정하지 않는다.}}

\section{세 점수체계와 해석 단위}
회사 입력형 비교기에서 서로 다른 세 수치를 같은 축으로 읽어서는 안 된다. 

\begingroup
\small
\setlength{\tabcolsep}{4pt}
\begin{longtable}{@{}L{3.2cm}L{2.8cm}L{8.3cm}@{}}
\caption{CPII 관련 세 점수체계의 구분}\label{tab:cpii_appendix_scale_distinction}\\
\toprule
수치 & 범위·부호 & 정확한 의미 \\
\midrule
\endfirsthead
\toprule
수치 & 범위·부호 & 정확한 의미 \\
\midrule
\endhead
회사별 CPII & $0$--$100$ & 고정 기준분포에서 feature별 부담 위치를 정해진 가중치로 합성한 회사별 구조적 부담점수. 기업 백분위가 아니다. \\
회사쌍 점수차 $\Delta\mathrm{CPII}_{A-B}$ & \(-100\)--\(+100\) & $\mathrm{CPII}(A)-\mathrm{CPII}(B)$이다. 양수는 A의 측정 부담 위치가 B보다 높고, 음수는 그 반대다. \\
legacy signed evidence score & \(-1\)--\(+1\) & 0을 중심으로 판단·증거를 코딩한 별도 지표다. feature 기반 CPII나 A-B 점수차가 아니며 같은 축·범례·해석 규칙을 쓰지 않는다. \\
\bottomrule
\end{longtable}
\endgroup

특히 $\Delta\mathrm{CPII}=0.10$점은 feature별 기준분포 위치 차이의 \textbf{가중합이 정확히 $0.001$}이라는 뜻이다($0.10/100=0.001$). 이는 비용 $0.1\%$, 매출 $0.1\%$, 회사 간 기업 백분위 $0.1\%$p를 뜻하지 않는다. 점수차는 먼저 feature 기여도로 분해한 뒤, 해당 원단위와 등록된 경제 노출기준이 있을 때만 채널별로 해석한다.

\section{고정 앵커와 운용 사양}
방법론 버전 \texttt{CPII\_FIRM\_INPUT\_EQW\_V1}은 고정 앵커 \texttt{POOL\_SCQ\_P01\_P99\_EMP\_RANK\_V1}을 사용한다. 입력 회사 집합을 추가·삭제·재정렬해도 기존 회사 점수는 바뀌지 않는다. 
현재 앵커에서 10개 활성 feature가 모두 관측된 결합 기준표본은 2/159행뿐이다. 저지원 활성 feature \texttt{cross\_market\_transfer\_index}, \texttt{foreign\_dependency\_ratio}의 합산 유효가중치는 22.22\%다. 따라서 회사 경험적 백분위, pair-gap 경험적 백분위, 산업 대표분포 주장은 비활성화한다.

운용 feature 구성은 $I_{\mathrm{reg}}$/$I_{\mathrm{risk}}$/$I_{\mathrm{market}}$ 순서로 3/4/3, 합계 10개다. 이 구성은 보고서의 더 넓은 이론 변수목록을 자동 대체하지 않는다.

\begin{align}
z_{ij} &= \mathcal{N}_{j}^{\mathrm{fixed}}\!\left(d_j x_{ij}\right), \\
I_{ik} &= \frac{100}{|J_k|}\sum_{j\in J_k} z_{ij}, \qquad k\in\{\mathrm{reg},\mathrm{risk},\mathrm{market}\}, \\
\mathrm{CPII}_i &= \frac{1}{3}\sum_k I_{ik}, \\
\Delta\mathrm{CPII}_{A-B} &= \mathrm{CPII}_A-\mathrm{CPII}_B =\sum_j 100w_j\left(z_{Aj}-z_{Bj}\right).
\end{align}
여기서 $d_j$는 부담 증가 방향을 맞추는 부호, $\mathcal{N}_{j}^{\mathrm{fixed}}$는 고정 앵커의 변환·1\%/99\% winsorization·경험순위 정규화, $w_j$는 하위지수 내부 동일가중과 세 축 동일가중을 전개한 유효 feature 가중치다.

\section{합성 회사 입력과 \texorpdfstring{$n$}{n}개 회사 비교}
\textbf{방법론 설명용 합성자료이며 실제 기업에 대한 추정이 아님.} 아래 네 회사는 원천 CSV에서 생성한 통제 합성 예시다. 회사명은 식별자일 뿐이며 값을 조회·추론·대체하지 않는다.

\begin{landscape}
\begingroup
\scriptsize
\setlength{\tabcolsep}{2.5pt}
\begin{longtable}{@{}L{1.2cm}L{4.1cm}L{1.7cm}L{1.0cm}L{2.1cm}L{2.1cm}L{2.1cm}L{2.1cm}@{}}
\caption{방법론 설명용 합성 회사의 10개 원단위 입력}\label{tab:cpii_appendix_synthetic_inputs}\\
\toprule
축 & 입력 feature & 단위 & 방향 & \WrapCell{\texttt{SYN-A}\\가온 콘텐츠} & \WrapCell{\texttt{SYN-B}\\누리 미디어} & \WrapCell{\texttt{SYN-C}\\다온 스토리} & \WrapCell{\texttt{SYN-D}\\라온 플레이} \\
\midrule
\endfirsthead
\toprule
축 & 입력 feature & 단위 & 방향 & \WrapCell{\texttt{SYN-A}\\가온 콘텐츠} & \WrapCell{\texttt{SYN-B}\\누리 미디어} & \WrapCell{\texttt{SYN-C}\\다온 스토리} & \WrapCell{\texttt{SYN-D}\\라온 플레이} \\
\midrule
\endhead
$I_{\mathrm{reg}}$ & \WrapCell{유효세율 격차\\\texttt{effective\_\allowbreak{}tax\_\allowbreak{}rate\_\allowbreak{}gap}} & 소수비율 & +1 & -0.0200 & -0.0550 & 0.0100 & -0.0050 \\
$I_{\mathrm{reg}}$ & \WrapCell{인프라 비용·해외의존 결합비율\\\texttt{infra\_\allowbreak{}dependency\_\allowbreak{}ratio}} & 소수비율 & +1 & 0.1332 & 0.0640 & 0.0986 & 0.1136 \\
$I_{\mathrm{reg}}$ & \WrapCell{네트워크 비용률\\\texttt{network\_\allowbreak{}cost\_\allowbreak{}ratio}} & 소수비율 & +1 & 0.1800 & 0.1050 & 0.1450 & 0.1600 \\
$I_{\mathrm{risk}}$ & \WrapCell{세전이익 기반 회수 대리값\\\texttt{cashflow\_\allowbreak{}recovery\_\allowbreak{}ratio}} & 배수 & -1 & 0.42 & 0.73 & 0.18 & 0.51 \\
$I_{\mathrm{risk}}$ & \WrapCell{자금조달비용 스프레드\\\texttt{funding\_\allowbreak{}cost\_\allowbreak{}spread}} & 소수비율 & +1 & 0.0120 & 0.0040 & 0.0180 & 0.0090 \\
$I_{\mathrm{risk}}$ & \WrapCell{투자 회수율\\\texttt{investment\_\allowbreak{}recovery\_\allowbreak{}ratio}} & 배수 & -1 & 3.20 & 4.10 & 2.65 & 3.55 \\
$I_{\mathrm{risk}}$ & \WrapCell{IP 자산화 결손\\\texttt{ip\_\allowbreak{}assetization\_\allowbreak{}gap}} & 소수비율 & +1 & 0.0800 & 0.0250 & 0.1100 & 0.0650 \\
$I_{\mathrm{market}}$ & \WrapCell{시장간 전이 대리지수\\\texttt{cross\_\allowbreak{}market\_\allowbreak{}transfer\_\allowbreak{}index}} & 0--1 지수 & +1 & 0.2800 & 0.1600 & 0.3400 & 0.2400 \\
$I_{\mathrm{market}}$ & \WrapCell{해외의존비율\\\texttt{foreign\_\allowbreak{}dependency\_\allowbreak{}ratio}} & 소수비율 & +1 & 0.7400 & 0.6100 & 0.6800 & 0.7100 \\
$I_{\mathrm{market}}$ & \WrapCell{비용 대비 설비투자 압박률\\\texttt{outflow\_\allowbreak{}to\_\allowbreak{}capex\_\allowbreak{}ratio}} & 배수 & +1 & 8.50 & 4.00 & 6.80 & 7.40 \\
\bottomrule
\end{longtable}
\endgroup
\end{landscape}

\begingroup
\small
\setlength{\tabcolsep}{4pt}
\begin{longtable}{@{}L{1.0cm}L{3.7cm}L{1.5cm}L{1.5cm}L{1.5cm}L{1.5cm}L{2.4cm}@{}}
\caption{합성 회사의 고정 앵커 점수와 부담순위}\label{tab:cpii_appendix_synthetic_ranking}\\
\toprule
순위 & 회사 & CPII & $I_{\mathrm{reg}}$ & $I_{\mathrm{risk}}$ & $I_{\mathrm{market}}$ & 관측/활성 feature \\
\midrule
\endfirsthead
\toprule
순위 & 회사 & CPII & $I_{\mathrm{reg}}$ & $I_{\mathrm{risk}}$ & $I_{\mathrm{market}}$ & 관측/활성 feature \\
\midrule
\endhead
1 & \WrapCell{\texttt{SYN-C}\\다온 스토리\\\texttt{foreign}} & 42.77 & 24.72 & 76.38 & 27.20 & 10/10 \\
2 & \WrapCell{\texttt{SYN-A}\\가온 콘텐츠\\\texttt{domestic}} & 39.90 & 12.74 & 73.36 & 33.59 & 10/10 \\
3 & \WrapCell{\texttt{SYN-D}\\라온 플레이\\\texttt{domestic}} & 35.98 & 11.02 & 69.38 & 27.55 & 10/10 \\
4 & \WrapCell{\texttt{SYN-B}\\누리 미디어\\\texttt{foreign}} & 32.22 & 5.09 & 65.06 & 26.50 & 10/10 \\
\bottomrule
\end{longtable}
\endgroup

\begingroup
\small
\setlength{\tabcolsep}{5pt}
\begin{longtable}{@{}L{3.0cm}L{2.0cm}L{2.0cm}L{2.0cm}L{2.0cm}@{}}
\caption{합성 회사 CPII 점수차 행렬(행 A $-$ 열 B, 점)}\label{tab:cpii_appendix_pair_matrix}\\
\toprule
행 $A$ $-$ 열 $B$ & \texttt{SYN-A} & \texttt{SYN-B} & \texttt{SYN-C} & \texttt{SYN-D} \\
\midrule
\endfirsthead
\toprule
행 $A$ $-$ 열 $B$ & \texttt{SYN-A} & \texttt{SYN-B} & \texttt{SYN-C} & \texttt{SYN-D} \\
\midrule
\endhead
\WrapCell{\texttt{SYN-A}\\가온 콘텐츠} & +0.00 & +7.68 & -2.87 & +3.91 \\
\WrapCell{\texttt{SYN-B}\\누리 미디어} & -7.68 & +0.00 & -10.55 & -3.77 \\
\WrapCell{\texttt{SYN-C}\\다온 스토리} & +2.87 & +10.55 & +0.00 & +6.78 \\
\WrapCell{\texttt{SYN-D}\\라온 플레이} & -3.91 & +3.77 & -6.78 & +0.00 \\
\bottomrule
\end{longtable}
\endgroup
표의 대각선은 0이고, 각 비대각 원소는 반대 방향 원소와 부호가 반대다. 순위와 행렬은 동일한 고정 앵커와 공통 10개 feature로 계산된다.

\section{2개 회사 상세 추적: SYN-A 대 SYN-B}
\texttt{SYN-A}와 \texttt{SYN-B}는 동일 cohort에서 비교 가능하며, 기준 점수차는 $\Delta\mathrm{CPII}_{A-B}=+7.68$점이다. 다음 표는 원단위 입력에서 고정 앵커 위치, 유효가중치, CPII 기여도까지 한 줄로 추적한다.

\begin{landscape}
\begingroup
\scriptsize
\setlength{\tabcolsep}{2.2pt}
\begin{longtable}{@{}L{0.8cm}L{3.9cm}L{1.4cm}L{1.4cm}L{2.0cm}L{1.5cm}L{1.5cm}L{1.5cm}L{1.8cm}@{}}
\caption{SYN-A 대 SYN-B의 feature별 점수차 분해}\label{tab:cpii_appendix_driver_trace}\\
\toprule
순위 & driver & 원값 A & 원값 B & 방향정렬 원차 & $z_A$ & $z_B$ & $w_j$ & 기여도(점) \\
\midrule
\endfirsthead
\toprule
순위 & driver & 원값 A & 원값 B & 방향정렬 원차 & $z_A$ & $z_B$ & $w_j$ & 기여도(점) \\
\midrule
\endhead
1 & \WrapCell{유효세율 격차\\\texttt{effective\_\allowbreak{}tax\_\allowbreak{}rate\_\allowbreak{}gap}\\$I_{\mathrm{reg}}$} & -0.0200 & -0.0550 & 0.0350 & 0.2534 & 0.1027 & 11.11\% & +1.67 \\
2 & \WrapCell{세전이익 기반 회수 대리값\\\texttt{cashflow\_\allowbreak{}recovery\_\allowbreak{}ratio}\\$I_{\mathrm{risk}}$} & 0.42 & 0.73 & 0.31 & 0.6755 & 0.4840 & 8.33\% & +1.60 \\
3 & \WrapCell{해외의존비율\\\texttt{foreign\_\allowbreak{}dependency\_\allowbreak{}ratio}\\$I_{\mathrm{market}}$} & 0.7400 & 0.6100 & 0.1300 & 0.2108 & 0.1024 & 11.11\% & +1.20 \\
4 & \WrapCell{자금조달비용 스프레드\\\texttt{funding\_\allowbreak{}cost\_\allowbreak{}spread}\\$I_{\mathrm{risk}}$} & 0.0120 & 0.0040 & 0.0080 & 0.7982 & 0.6579 & 8.33\% & +1.17 \\
5 & \WrapCell{비용 대비 설비투자 압박률\\\texttt{outflow\_\allowbreak{}to\_\allowbreak{}capex\_\allowbreak{}ratio}\\$I_{\mathrm{market}}$} & 8.50 & 4.00 & 4.50 & 0.2969 & 0.1927 & 11.11\% & +1.16 \\
6 & \WrapCell{인프라 비용·해외의존 결합비율\\\texttt{infra\_\allowbreak{}dependency\_\allowbreak{}ratio}\\$I_{\mathrm{reg}}$} & 0.1332 & 0.0640 & 0.0692 & 0.0714 & 0.0238 & 11.11\% & +0.53 \\
7 & \WrapCell{네트워크 비용률\\\texttt{network\_\allowbreak{}cost\_\allowbreak{}ratio}\\$I_{\mathrm{reg}}$} & 0.1800 & 0.1050 & 0.0750 & 0.0573 & 0.0260 & 11.11\% & +0.35 \\
8 & \WrapCell{시장간 전이 대리지수\\\texttt{cross\_\allowbreak{}market\_\allowbreak{}transfer\_\allowbreak{}index}\\$I_{\mathrm{market}}$} & 0.2800 & 0.1600 & 0.1200 & 0.5000 & 0.5000 & 11.11\% & +0.00 \\
9 & \WrapCell{투자 회수율\\\texttt{investment\_\allowbreak{}recovery\_\allowbreak{}ratio}\\$I_{\mathrm{risk}}$} & 3.20 & 4.10 & 0.90 & 0.6968 & 0.6968 & 8.33\% & +0.00 \\
10 & \WrapCell{IP 자산화 결손\\\texttt{ip\_\allowbreak{}assetization\_\allowbreak{}gap}\\$I_{\mathrm{risk}}$} & 0.0800 & 0.0250 & 0.0550 & 0.7637 & 0.7637 & 8.33\% & +0.00 \\
\midrule
\multicolumn{8}{r}{기여도 합계} & +7.678162 \\
\multicolumn{8}{r}{합계 $-$ 표시 전 원점수차 잔차} & -5.33e-15 \\
\bottomrule
\end{longtable}
\endgroup
\end{landscape}

절대기여도 상위 3개 driver는 유효세율 격차 +1.67점, 세전이익 기반 회수 대리값 +1.60점, 해외의존비율 +1.20점이다. 양의 기여는 A-B 점수차를 넓히고 음의 기여는 줄이며, 0 기여도도 원값 차이가 없다는 뜻은 아닐 수 있다. 고정 경험순위에서 같은 위치로 매핑되면 원값 차이가 있어도 점수 기여는 0이 된다.

\subsection{원단위 차이에서 경제량까지}
공통규모 상당액은 두 회사에 공통으로 적용한 노출기준에서 driver의 원단위 차이를 비교한다. 실제규모 차이는 회사별 수준과 회사별 노출액을 각각 결합하므로 공통규모 값과 다른 질문에 답한다. \texttt{scenario\_only}는 관측 손익이 아니라 명시적 가정에 따른 시나리오이며, \texttt{non\_monetizable}은 독립적·비중복 경제계수가 식별되지 않아 금액으로 바꾸지 않는다.

\begin{landscape}
\begingroup
\scriptsize
\setlength{\tabcolsep}{2pt}
\begin{longtable}{@{}L{4.3cm}L{3.1cm}L{3.1cm}L{4.1cm}L{4.1cm}L{2.5cm}@{}}
\caption{SYN-A 대 SYN-B의 경제해석 low/base/high(2025년 KRW, 억원)}\label{tab:cpii_appendix_economic_trace}\\
\toprule
driver & 해석 유형 & 노출기준 & 공통규모 low / base / high & 실제규모 low / base / high & 중복군 \\
\midrule
\endfirsthead
\toprule
driver & 해석 유형 & 노출기준 & 공통규모 low / base / high & 실제규모 low / base / high & 중복군 \\
\midrule
\endhead
\WrapCell{유효세율 격차\\\texttt{effective\_\allowbreak{}tax\_\allowbreak{}rate\_\allowbreak{}gap}} & 조건부 회계 상당액 & 양(+) 세전이익 & +34.65 / +38.50 / +42.35 & +54.90 / +61.00 / +67.10 & 조세부담 \\
\WrapCell{세전이익 기반 회수 대리값\\\texttt{cashflow\_\allowbreak{}recovery\_\allowbreak{}ratio}} & 조건부 성과 상당액 & 콘텐츠투자액 & +558.00 / +620.00 / +682.00 & +351.00 / +390.00 / +429.00 & 회수성 \\
\WrapCell{해외의존비율\\\texttt{foreign\_\allowbreak{}dependency\_\allowbreak{}ratio}} & 화폐환산 불가 & 적용 없음 & -- & -- & 네트워크·인프라 \\
\WrapCell{자금조달비용 스프레드\\\texttt{funding\_\allowbreak{}cost\_\allowbreak{}spread}} & 조건부 회계 상당액 & 평균 이자부부채 & +18.72 / +20.80 / +22.88 & +24.48 / +27.20 / +29.92 & 자금·IP \\
\WrapCell{비용 대비 설비투자 압박률\\\texttt{outflow\_\allowbreak{}to\_\allowbreak{}capex\_\allowbreak{}ratio}} & 화폐환산 불가 & 적용 없음 & -- & -- & 네트워크·인프라 \\
\WrapCell{인프라 비용·해외의존 결합비율\\\texttt{infra\_\allowbreak{}dependency\_\allowbreak{}ratio}} & 화폐환산 불가 & 적용 없음 & -- & -- & 네트워크·인프라 \\
\WrapCell{네트워크 비용률\\\texttt{network\_\allowbreak{}cost\_\allowbreak{}ratio}} & 조건부 회계 상당액 & 매출 & +810.00 / +900.00 / +990.00 & +297.00 / +330.00 / +363.00 & 네트워크·인프라 \\
\WrapCell{시장간 전이 대리지수\\\texttt{cross\_\allowbreak{}market\_\allowbreak{}transfer\_\allowbreak{}index}} & 가정 기반 시나리오 & 대상시장 매출 & +8.42 / +11.70 / +15.44 & +10.08 / +14.00 / +18.48 & 시장전이 \\
\WrapCell{투자 회수율\\\texttt{investment\_\allowbreak{}recovery\_\allowbreak{}ratio}} & 조건부 성과 상당액 & 콘텐츠투자액 & +1,620.00 / +1,800.00 / +1,980.00 & +306.00 / +340.00 / +374.00 & 회수성 \\
\WrapCell{IP 자산화 결손\\\texttt{ip\_\allowbreak{}assetization\_\allowbreak{}gap}} & 가정 기반 시나리오 & 매출 & +594.00 / +660.00 / +726.00 & +405.00 / +450.00 / +495.00 & 자금·IP \\
\bottomrule
\end{longtable}
\endgroup
\end{landscape}
모든 경제해석 행의 집계 규칙은 \texttt{never\_auto\_sum}이다. 조건부 회계 상당액, 조건부 성과 상당액, 시나리오 금액은 개념과 중복 family가 다르고 서로 원인·결과 관계일 수도 있으므로 합계 손실액처럼 자동 합산하지 않는다. low/base/high도 결정론적 입력·노출 가정 범위이며 신뢰구간이나 확률구간이 아니다.

\section{같은 점수차, 다른 경제적 크기: 골든 통제 예시}
두 골든 pair는 하나의 driver만 바꾸어 동일한 CPII 점수차를 만들지만 노출기준과 원단위가 다르다.

\begingroup
\small
\setlength{\tabcolsep}{4pt}
\begin{longtable}{@{}L{3.2cm}L{3.0cm}L{2.2cm}L{2.8cm}L{2.8cm}@{}}
\caption{동일 CPII 점수차의 네트워크·조달금리 경제량 대비}\label{tab:cpii_appendix_equal_gap_golden}\\
\toprule
통제 pair & 단일 변경 driver & $\Delta$CPII(점) & 공통 노출 base(억원) & 공통규모 base(억원) \\
\midrule
\endfirsthead
\toprule
통제 pair & 단일 변경 driver & $\Delta$CPII(점) & 공통 노출 base(억원) & 공통규모 base(억원) \\
\midrule
\endhead
\WrapCell{네트워크 비용률 통제\\\texttt{GOLD-N-A}\\\texttt{GOLD-N-B}} & 네트워크 비용률 & +1.3889 & 10,000.00 & 997.52 \\
\WrapCell{조달금리 통제\\\texttt{GOLD-F-A}\\\texttt{GOLD-F-B}} & 조달금리 스프레드 & +1.3889 & 1,000.00 & 2.00 \\
\bottomrule
\end{longtable}
\endgroup
두 pair의 점수차는 각각 1.3889점으로 같지만, 네트워크 비용률 채널의 공통규모 base는 약 997.52억원, 자금조달비용 스프레드 채널은 약 2.00억원이다. 전자는 후자의 약 498.95배다. 따라서 보편적인 ``CPII 1점당 원화'' 계수는 존재하지 않는다.

\section{민감도 판독과 한계}
\begingroup
\footnotesize
\setlength{\tabcolsep}{2pt}
\begin{longtable}{@{}L{2.5cm}L{1.5cm}L{2.8cm}L{1.5cm}L{2.0cm}L{3.8cm}@{}}
\caption{합성 회사 pair별 결정론적 민감도 범위}\label{tab:cpii_appendix_sensitivity}\\
\toprule
pair & 기준 gap & 민감도 구간 [min, max] & 부호 안정 & 양(+) 사양 비중 & 판독 \\
\midrule
\endfirsthead
\toprule
pair & 기준 gap & 민감도 구간 [min, max] & 부호 안정 & 양(+) 사양 비중 & 판독 \\
\midrule
\endhead
SYN-A–SYN-B & +7.68 & [-2.00, +14.32] & \texttt{no} & 83.33\% & 차이 불명확 \\
SYN-A–SYN-C & -2.87 & [-11.47, +4.39] & \texttt{no} & 16.67\% & 차이 불명확 \\
SYN-A–SYN-D & +3.91 & [-7.21, +11.14] & \texttt{no} & 83.33\% & 차이 불명확 \\
SYN-B–SYN-C & -10.55 & [-15.88, -3.52] & \texttt{yes} & 0.00\% & 시험 사양 내 방향 유지 \\
SYN-B–SYN-D & -3.77 & [-11.61, +3.22] & \texttt{no} & 16.67\% & 차이 불명확 \\
SYN-C–SYN-D & +6.78 & [-1.69, +12.69] & \texttt{no} & 83.33\% & 차이 불명확 \\
\bottomrule
\end{longtable}
\endgroup
선택 pair SYN-A–SYN-B의 범위는 -2.00점에서 +14.32점이며 0을 포함한다. 따라서 기준점수의 부호만으로 안정된 차이라고 결론내리지 않고 차이 불명확로 판독한다.

\begingroup\small
\begin{itemize}[leftmargin=2em,itemsep=0.15em,topsep=0.25em]
\item 민감도 [min, max] 구간은 입력 범위와 사전 정의 대체 사양을 통과시킨 결정론적 범위다. 신뢰구간·확률구간·예측구간이 아니며, 양(+) 사양 비중도 확률이 아니다.
\item 결합 기준표본이 얇으므로 회사 경험적 백분위, pair-gap 경험적 백분위, 산업 대표성은 주장하지 않는다. 저지원 feature의 가중치가 결과에 포함된다는 Research Beta 경고를 항상 함께 본다.
\item 합성 예시는 계산 재현성을 검증할 뿐 실제 회사, 산업 평균, 법적 차별·역차별, 위법성, 손해 또는 인과효과의 증거가 아니다.
\end{itemize}
\endgroup

% End of auto-generated CPII company-comparator appendix.
